\chapter{Experimento}
\label{cap:experimento}

\section{Descrição do Experimento}
\label{sec:descricao}

O experimento é a implementação do \textbf{IBE Mail}: uma aplicação de e-mail com interface
gráfica que demonstra o esquema FullIdent de ponta a ponta. A aplicação roda com múltiplas
identidades na mesma máquina e tem um PKG embutido para registro, autenticação e extração de
chaves. Não há servidor separado, o próprio PKG faz parte da própria aplicação, por razões de
simplificação.

\subsection{O que demonstramos}

O objetivo é mostrar, na prática, que o IBE funciona como descrito:

\begin{enumerate}
  \item Cifrar para \texttt{alice@ufsc.br} sem precisar de nenhuma chave pública de Alice
    antecipadamente.

  \item Ver a decifragem falhar com a chave errada, com o erro visível na interface.

  \item Revogar uma identidade no PKG, bloqueando novas extrações de chave para ela.

  \item Cifrar para \texttt{alice@ufsc.br|2026-06}, de forma que Alice só consegue decifrar
    se tiver obtido a chave do período correspondente.
\end{enumerate}


\subsection{Formato de Mensagem (\texttt{.ibemail})}

Cada mensagem é armazenada em disco como arquivo texto com cabeçalhos e blocos cifrados em
base64:

\begin{verbatim}
IBE-MAIL/1.0
Message-ID: 550e8400-e29b-41d4-a716-446655440000
From: alice@ufsc.br
To: bob@ufsc.br
Date: 2026-06-21T14:30:22
---
U: <base64 do ponto G1>
V: <base64 de n bytes>
W: <base64 de n bytes>
\end{verbatim}

Assunto e corpo vão cifrados dentro da mensagem, sem aparecer no header.

\subsection{Registro de Identidades (\texttt{registry.dat})}

O PKG mantém um arquivo com as identidades cadastradas, suas senhas em SHA-256 e se foram
aprovadas:

\begin{verbatim}
alice@ufsc.br   a665a...ae3   approved
bob@ufsc.br     5e884...d8    approved
carol@ufsc.br   <hash>        pending
\end{verbatim}

Um novo usuário se registra com identidade e senha; o PKG aprova o cadastro pelo painel de
gerenciamento. Apenas usuários aprovados podem obter sua chave privada e logar no sistema.

% ---

\section{Pseudocódigo}
\label{sec:pseudocodigo}

Os pseudocódigos a seguir descrevem as operações principais da aplicação.

\subsection{Inicialização do Sistema (Setup)}

\begin{algorithm}[H]
\caption{IBE Mail -- Setup}
\begin{algorithmic}[1]
\Procedure{Setup}{}
  \State Carregar parâmetros Tipo A da biblioteca PBC (\texttt{a.param})
  \State Inicializar grupo $G_1$ de ordem $q$ com gerador $P$
  \State Escolher $s \xleftarrow{\$} \mathbb{Z}_q^*$ \Comment{segredo mestre}
  \State Calcular $P_{\mathit{pub}} \leftarrow s \cdot P$
  \State Serializar e salvar \textit{params} em \texttt{data/params.pub} (base64)
  \State Serializar e salvar $s$ em \texttt{data/master.sec} (base64)
  \State Criar diretórios \texttt{data/keys/}, \texttt{data/mail/inbox/}, \texttt{data/mail/sent/}
  \State Inicializar \texttt{data/registry.dat} vazio
\EndProcedure
\end{algorithmic}
\end{algorithm}

\subsection{Extração de Chave Privada (Extract)}

\begin{algorithm}[H]
\caption{IBE Mail -- Extract}
\begin{algorithmic}[1]
\Procedure{Extract}{$\mathit{ID}$, $\mathit{password}$}
  \State Carregar registro de \texttt{data/registry.dat}
  \If{$\mathit{ID}$ não está no registro ou status $\neq$ \textit{approved}}
    \Return $\perp$
  \EndIf
  \If{$\mathit{SHA256}(\mathit{password}) \neq \text{hash armazenado}$}
    \Return $\perp$ \Comment{autenticação falhou}
  \EndIf
  \State Carregar $s$ de \texttt{data/master.sec}
  \State Carregar \textit{params} de \texttt{data/params.pub}
  \State $Q_{\mathit{ID}} \leftarrow H_1(\mathit{ID})$ \Comment{\texttt{element\_from\_hash} da PBC}
  \State $d_{\mathit{ID}} \leftarrow s \cdot Q_{\mathit{ID}}$
  \State Serializar e salvar $d_{\mathit{ID}}$ em \texttt{data/keys/}\textit{ID}\texttt{.prv}
  \Return $d_{\mathit{ID}}$
\EndProcedure
\end{algorithmic}
\end{algorithm}

\subsection{Cifragem de Mensagem (Encrypt)}

\begin{algorithm}[H]
\caption{IBE Mail -- Encrypt}
\begin{algorithmic}[1]
\Procedure{Encrypt}{$\mathit{ID}_{\mathit{dest}}$, $M$}
  \State Carregar \textit{params} de \texttt{data/params.pub}
  \State $Q_{\mathit{ID}} \leftarrow H_1(\mathit{ID}_{\mathit{dest}})$
  \State Escolher $\sigma \xleftarrow{\$} \{0,1\}^n$
  \State $r \leftarrow H_3(\sigma, M)$
  \State $U \leftarrow r \cdot P$
  \State $V \leftarrow \sigma \oplus H_2\!\left(\hat{e}(Q_{\mathit{ID}},\, P_{\mathit{pub}})^r\right)$
  \State $W \leftarrow M \oplus H_4(\sigma)$
  \State $\mathit{uuid} \leftarrow \text{gerar UUID v4}$
  \State Escrever arquivo \texttt{data/mail/inbox/}\textit{ID\_dest}\texttt{/}\textit{uuid}\texttt{.ibemail}
  \State Escrever cópia em \texttt{data/mail/sent/}\textit{ID\_orig}\texttt{/}\textit{uuid}\texttt{.ibemail}
  \Return $\langle U, V, W \rangle$
\EndProcedure
\end{algorithmic}
\end{algorithm}

\subsection{Decifração de Mensagem (Decrypt)}

\begin{algorithm}[H]
\caption{IBE Mail -- Decrypt}
\begin{algorithmic}[1]
\Procedure{Decrypt}{$\langle U, V, W \rangle$, $d_{\mathit{ID}}$}
  \If{$d_{\mathit{ID}}$ não está carregada}
    \Return $\perp$ \Comment{usuário não possui chave privada}
  \EndIf
  \State $\sigma' \leftarrow V \oplus H_2(\hat{e}(d_{\mathit{ID}},\, U))$
  \State $M' \leftarrow W \oplus H_4(\sigma')$
  \State $r' \leftarrow H_3(\sigma', M')$
  \If{$U \neq r' \cdot P$}
    \Return $\perp$ \Comment{verificação falhou: chave incorreta ou texto cifrado inválido}
  \EndIf
  \Return $M'$
\EndProcedure
\end{algorithmic}
\end{algorithm}

\subsection{Fluxo de Uso da Interface}

O uso da aplicação segue os passos abaixo:

\begin{enumerate}
  \item \textbf{Registro.} O usuário preenche identidade e senha no forms de registro.
    O PKG registra o usuário como pendente.

  \item \textbf{Aprovação.} O administrador aprova o cadastro no painel do PKG, mudando
    o status para \textit{approved}.

  \item \textbf{Obtenção de chave.} O usuário informa sua senha no painel esquerdo. A
    aplicação executa o Extract e salva $d_{\mathit{ID}}$ em disco; a chave fica
    disponível na sessão.

  \item \textbf{Envio de mensagem.} Pelo modal de envio de mensagem, o remetente informa a 
      identidade do destinatário, assunto e corpo. A aplicação executa o Encrypt e o arquivo
    \texttt{.ibemail} vai para o inbox do destinatário.

  \item \textbf{Leitura.} Ao selecionar uma mensagem no painel central, a aplicação executa o
    Decrypt com a chave do usuário ativo. O texto decifrado aparece no painel direito, com um
    indicador \textit{Decifrado com alice@ufsc.br}. Se a decifragem falha, a interface mostra
    um erro.

  \item \textbf{Revogação.} O administrador revoga uma identidade no painel do PKG. O status
    muda para \textit{revoked}, o PKG aapaga a chave do disco e bloqueia novas extrações.
    Mensagens cifradas para identidades com período expirado ficam inacessíveis, já que não é
    mais possível obter a chave.
\end{enumerate}
